\section{Proofs}
\label{sec:proofs}
Throughout, the grid has $N$ cells with periodic indexing, and $\operatorname{roll}_a$
shifts a field by one cell along axis $a$ with wrap-around, so it is a permutation of the
cells.

\paragraph{Proposition~\ref{prop:conservation}.} Summing Equation~\ref{eq:flux} over all
cells,
\begin{equation}
 \sum_{\mathrm{cells}}s_{t+1}-\sum_{\mathrm{cells}}s_t
 =\sum_{a=1}^{d}\Bigl(\sum_{\mathrm{cells}}\operatorname{roll}_a(F_a)
 -\sum_{\mathrm{cells}}F_a\Bigr)=0,
\end{equation}
because a sum over all cells is invariant under a permutation of the cells. Nothing in
the argument depends on how $F_a$ was computed, so the identity holds for every parameter
value, every state and every perception stencil, including the dilated perception of
PI-NCA-MS. In floating point each sum is exact up to round-off, which is what the
measured mass errors of order $10^{-7}$ relative reflect.

\paragraph{Proposition~\ref{prop:perfield}.} Equation~\ref{eq:mc} applies the scalar
update to each channel $c$ with its own flux $F^{(c)}$. Applying the argument above to
channel $c$ alone gives the result for that channel; the dependence of $F^{(c)}$ on the
other channels plays no role. By contrast, a single lumped constraint
$\sum_c\sum_{\mathrm{cells}}s^{(c)}$ admits updates with
$\sum_{\mathrm{cells}}\Delta s^{(1)}=-\sum_{\mathrm{cells}}\Delta s^{(2)}\neq0$, which
violate the conservation of each field.

\paragraph{Proposition~\ref{prop:headroom}.} Take $d=M_t-\sum_j\widetilde s_j>0$; the case
$d<0$ is symmetric with $r_i=\widetilde s_i-\ell$, and $d=0$ leaves the state unchanged.
Here $r_i=h-\widetilde s_i\ge0$ because $\widetilde s_i\le h$, and
$C_r=\sum_jr_j=Nh-\sum_j\widetilde s_j$. Let $\lambda=\min(d/C_r,1)\in(0,1]$. Each output
is $\widetilde s_i+\lambda(h-\widetilde s_i)$, a convex combination of $\widetilde s_i$
and $h$, and so lies in $[\widetilde s_i,h]\subseteq[\ell,h]$. The total added is
$\lambda C_r$. If $M_t\le Nh$ then $d\le C_r$, so $\lambda=d/C_r$ and the total added is
exactly $d$, giving $\sum\mathcal{R}_{M_t}(\widetilde s)=M_t$. If $M_t>Nh$ then
$\lambda=1$ and every cell equals $h$, the nearest admissible state. If $C_r=0$ every cell
is already at $h$ and the update is taken to be the identity. The map is a composition of
elementary operations with a single $\min$, so it is differentiable except on the set
$|d|=C_r$. For the uniform offset, by contrast, a cell with $\widetilde s_i=h$ receives
$h+d/N>h$ whenever $d>0$, so that restoration cannot guarantee the bound.

\section{PDE Definitions and Teacher Scope}
\label{sec:equations}
The experiments concern the numerical one-step map for each PDE. The governing
forms below define the structural distinctions used in the main text; they do not
replace the full teacher configuration. Domain sizes, coefficients, solver
substeps, initial-condition distributions, and numerical precision must be taken
from the corresponding experiment configuration to reproduce a particular score.

\paragraph{Heat.}
\begin{equation}
 \partial_t u=\alpha\nabla^2u.
\end{equation}
With periodic boundaries and constant diffusivity, spatial integration removes the
Laplacian and conserves total heat. This is the smooth local-diffusion test.

\paragraph{Advection--diffusion.}
\begin{equation}
 \partial_t u=-\mathbf{c}\cdot\nabla u+\alpha\nabla^2u.
\end{equation}
Constant-velocity periodic transport conserves the integral as well as moving the
field. Its FNO result should not be interpreted as evidence that conservation is
absent; it shows that a correct invariant alone need not determine the best model.

\paragraph{Allen--Cahn.}
\begin{equation}
 \partial_t u=\varepsilon^2\nabla^2u-W'(u),\qquad W'(u)=u^3-u.
\end{equation}
The reaction term generally changes the spatial integral. Evolving interfaces
provide a setting in which both non-conservation and spatial detail matter.

\paragraph{Wave.}
\begin{equation}
 \partial_{tt}u=c^2\nabla^2u.
\end{equation}
The emulator represents this as a first-order system in $(u,\partial_tu)$.
Propagation and coupling between fields distinguish it from scalar diffusion.

\paragraph{Shallow-water.}
The source-free periodic shallow-water system is represented in conserved variables
$(h,hu,hv)$. Per-field conservation concerns height and momentum, not arbitrary
primitive velocity channels. This variable choice is part of the architectural prior.

\paragraph{Nagumo and FitzHugh--Nagumo.}
The scalar bistable Nagumo equation is
\begin{equation}
 \partial_tu=D\nabla^2u+u(1-u)(u-a).
\end{equation}
FHN is the two-field activator--inhibitor reaction system in the implementation.
Its reaction coefficients and precise variant are not specified by the supplied
summary. For both problems, reaction terms generally preclude conserving each
field's total independently.

\paragraph{Cahn--Hilliard.}
\begin{equation}
 \partial_tu=\nabla^2\bigl(-\varepsilon^2\nabla^2u+W'(u)\bigr),
 \qquad W'(u)=u^3-u.
\end{equation}
This fourth-order equation conserves mass on a periodic domain. The reported
teacher clips to $[-1,1]$, and the emulator bounding ablation follows that
convention. We do not identify this clipped map with the unconstrained continuum
equation. The learned correction also does not establish a discrete free-energy
law. These distinctions are important when interpreting the stability improvement.

\paragraph{Gray--Scott.}
\begin{align}
 \partial_tu&=D_u\nabla^2u-uv^2+F(1-u),\\
 \partial_tv&=D_v\nabla^2v+uv^2-(F+k)v.
\end{align}
The feed and removal terms change field totals. The reference parameters with
$dt=2.0$ were reported to diverge; the benchmark uses $dt=1.0$. The reported
errors remain large, so this is a difficult case rather than an unqualified success
for either family.

\paragraph{Navier--Stokes.}
The reported vorticity formulation uses the Poisson relation
$\nabla^2\psi=-\omega$ to reconstruct the flow. This is the relevant source of
global coupling. Viscosity, forcing, the treatment of the zero Fourier mode, and
the full time-discretization settings must be specified by the experiment code;
no additional values are inferred here.

\paragraph{Three-dimensional suite.}
Heat, advection--diffusion, Allen--Cahn, Nagumo, Gray--Scott, and FHN are implemented
in NDHWC layout and evaluated on $24^3$ grids in this run. The implementation uses a
7-point Laplacian, flux divergence along all three axes, and spectral convolution
through \model{rfftn}; the explicit stability bound tightens to
$\Delta t\le\Delta x^2/(6D)$ in three dimensions, which is why Gray--Scott uses
$\Delta t=0.5$ here against $1.0$ in 2-D. The 3-D correctness subset contains 11 checks.
It does not include 3-D versions of the wave, shallow-water, Navier--Stokes, or
Cahn--Hilliard problems.

\section{Architecture Details and Schematics}
\label{sec:architecture-details}
The model names in the paper map to the released implementation as follows.
\textbf{PI-NCA} is \model{DeepFluxNCA} (\model{pi\_nca}): a $3\times3$ circular
perception convolution with 32 channels, a per-cell MLP with 64 and 32 hidden channels,
and a zero-initialized flux head with $2C$ outputs. \textbf{MC-PI-NCA} is
\model{MultiChannelFluxNCA} (\model{mc\_flux\_nca}), the same per-field construction with
48 perception and 96 hidden channels. \textbf{PI-NCA-B} is \model{BoundedConsFluxNCA}
(\model{bounded\_cons\_nca}), PI-NCA followed by the clip and headroom projection.
\textbf{PI-NCA-MS} is \model{MultiScaleFluxNCA} (\model{multiscale\_flux\_nca}) with
dilations $(1,2,4)$ and 24 features per dilation, and \textbf{PI-NCA-MS-B}
(\model{bounded\_multiscale\_nca}) adds the projection. \textbf{PI-NCA-Spec} is
\model{SpectralFluxNCA} (\model{spectral\_flux\_nca}), a local flux stream plus a
two-layer Fourier stream of width 16 with 8 modes, followed by mass restoration.
\textbf{NCA} is the standard residual NCA (\model{plain\_nca}); \textbf{FNO-S},
\textbf{ResNet-S} and \textbf{U-Net-S} are the parameter-matched controls. On single-field
problems PI-NCA and MC-PI-NCA differ only in width; the per-field divergence is what
distinguishes MC-PI-NCA from a lumped construction on multi-field problems, and the
compact PI-NCA configuration uses the same per-field divergence there.

For a fixed architecture width and stencil, local convolutional update cost scales
linearly with the number of cells. Exact FLOP counts also depend on channel widths,
MLP depth, and stencil support; parameter count alone is not a runtime model.
The update heads are reported to be zero-initialized to start residual emulator
maps at the identity. This initialization statement is not intended to describe
the native prediction functions of the continuous PINN or DeepONet.

The multi-scale update can be written as
\begin{equation}
 s_{t+1}=\mathcal{C}\!\left[s_t+D_hF_\theta
 \bigl([P_1(s_t)\Vert P_2(s_t)\Vert P_4(s_t)]\bigr)\right],
\end{equation}
where $\mathcal{C}$ is the identity when \model{bounds=None} and the reported
clipping/restoration operation when bounds are enabled. In the source discussion,
\model{bounds=(-1,1)} is reported to match the CH winner with approximately
$6\times$ lower variance; the supporting variance panel is not supplied, so this
remains a contextual observation rather than a headline result.

\paragraph{FNO.}
A spectral layer combines a pointwise transformation with learned Fourier weights:
\begin{equation}
 v_{\ell+1}=\sigma\!\left(Wv_\ell+
 \mathcal{F}^{-1}\bigl(R_\ell\odot\mathcal{F}[v_\ell]\bigr)\right).
\end{equation}
The implementation truncates spectral mixing to the lowest $m\times m$ modes in
2-D. Global mixing, mode truncation, and pointwise nonlinearities jointly determine
its behavior. The present experiments neither imply that FNO always removes sharp
features nor guarantee invariant accuracy across grids.

\paragraph{PINN.}
The continuous single-IVP baseline minimizes
\begin{equation}
\begin{split}
 \mathcal{L}_{\mathrm{PINN}}={}&\lambda_r\|\partial_tu_\phi-\mathcal{N}[u_\phi]\|^2
 +\lambda_b\|\mathcal{B}[u_\phi]\|^2\\
 &+\lambda_0\|u_\phi(\cdot,0)-u_0\|^2.
\end{split}
\end{equation}
It uses the equation, boundary conditions, and initial condition without labeled
interior solver pairs. Its continuous coordinate representation permits off-grid
queries, but the experiment trains anew for each IVP.

\paragraph{DeepONet.}
The branch network encodes the input function at sensors and the trunk encodes the
query coordinate, giving
$G(a)(y)\approx\sum_k b_k(a)t_k(y)$. It is trained on solver-generated function
pairs across initial conditions. Its supervised operator task differs from both the
single-IVP PINN and the fixed-step autoregressive emulator.

\begin{table}[t]
\centering
\caption{Structural properties of the PI-NCA family and the baselines. Perception extent
is the reach of a single update, not the full receptive field after the divergence.
Parameter counts are for the single-field heat configuration; on $C$-field problems the
PI-NCA flux head grows by $2C$ output channels only. ``Exact'' conservation holds for
any weights, at initialization and after training, up to floating-point round-off.}
\label{tab:axes}
\resizebox{\linewidth}{!}{\begin{tabular}{lcccccc}
\toprule
Property & NCA & PI-NCA & MC-PI-NCA & PI-NCA-MS(-B) & ResNet / U-Net & FNO \\
\midrule
Perception extent  & $\pm1$ cell & $\pm1$ cell & $\pm1$ cell & $\pm4$ cells & local / pooled & global \\
Mass conservation  & none & \textbf{exact} & \textbf{exact, per field} & \textbf{exact, per field} & none & none \\
Admissible bounds  & no & no & no & \textbf{with -B} & no & no \\
Params (heat)      & 6{,}784 & \textbf{4{,}576} & 9{,}936 & 5{,}520 & 74k / 265k & $5.9\!\times\!10^{5}$ \\
Shared local rule  & yes & yes & yes & yes & no & no \\
Boundary handling  & periodic & periodic & periodic & periodic & periodic (circ.) & periodic (FFT) \\
\bottomrule
\end{tabular}}
\end{table}


The schematics below retain the original architecture decompositions. Arrows
describe computation, not guarantees of numerical stability. In particular,
``restore mass'' denotes the reported correction routine; simultaneous bound and
mass preservation requires checking that routine, not merely composing two labels.
\begin{figure}[t]
\centering
\resizebox{\linewidth}{!}{\begin{tikzpicture}[node distance=6mm and 6mm]
  \node[box] (x) {state $s_t$\\$(B,H,W,C)$};
  \node[box, right=of x] (p) {perceive\\$3{\times}3$ circ., 32 ch};
  \node[box, right=of p] (m) {$1{\times}1$ MLP\\64, 32 ch, ReLU};
  \node[box, right=of m] (f) {flux head\\$2C$ ch, zero-init};
  \node[box, right=of f] (d) {discrete\\divergence};
  \node[op, right=6mm of d] (add) {$+$};
  \node[box, right=of add] (o) {$s_{t+1}$\\mass-conserving};
  \draw[flow] (x) -- (p);
  \draw[flow] (p) -- (m);
  \draw[flow] (m) -- (f);
  \draw[flow] (f) -- (d);
  \draw[flow] (d) -- (add);
  \draw[flow] (add) -- (o);
  \draw[flow] (x.south) to[out=-40,in=-140] (add.south);
\end{tikzpicture}}
\caption{PI-NCA (\texttt{DeepFluxNCA}). The head predicts a flux per axis; its periodic
discrete divergence sums to zero over the grid, so the residual update conserves mass by
construction.}
\label{fig:deepflux}
\end{figure}

\begin{figure}[t]
\centering
\resizebox{\linewidth}{!}{\begin{tikzpicture}[node distance=6mm and 6mm]
  \node[box] (x) {state $s_t$\\$(B,H,W,C)$};
  \node[box, right=of x] (p) {perceive\\48 ch\,$\to$\,MLP 96 ch};
  \node[box, right=of p] (f) {flux head\\$2C$ channels};
  \node[box, right=of f] (d) {per-channel\\divergence};
  \node[op, right=6mm of d] (add) {$+$};
  \node[box, right=of add] (o) {$s_{t+1}$\\each field's\\mass conserved};
  \draw[flow] (x) -- (p);
  \draw[flow] (p) -- (f);
  \draw[flow] (f) -- (d);
  \draw[flow] (d) -- (add);
  \draw[flow] (add) -- (o);
  \draw[flow] (x.south) to[out=-40,in=-140] (add.south);
\end{tikzpicture}}
\caption{MC-PI-NCA (\texttt{MultiChannelFluxNCA}). Perception sees all fields jointly;
the flux head emits a flux pair $(f_x,f_y)$ per field, and a per-channel divergence
conserves each field's total independently.}
\label{fig:multichannel}
\end{figure}

\begin{figure}[t]
\centering
\resizebox{\linewidth}{!}{\begin{tikzpicture}[node distance=6mm and 6mm]
  \node[box] (x) {state $s_t$\\record $M_t$};
  \node[box, right=of x] (u) {PI-NCA\\flux update};
  \node[box, right=of u] (cl) {clip\\$[\ell,h]$};
  \node[box, right=of cl] (ce) {headroom\\projection $\mathcal{R}_{M_t}$};
  \node[box, right=of ce] (o) {$s_{t+1}$\\bounded and\\mass-exact};
  \draw[flow] (x) -- (u);
  \draw[flow] (u) -- (cl);
  \draw[flow] (cl) -- (ce);
  \draw[flow] (ce) -- (o);
  \draw[flow] (x.south) -- ++(0,-8mm) -| (ce.south);
\end{tikzpicture}}
\caption{PI-NCA with the bound projection (\texttt{BoundedConsFluxNCA}). The clip enforces
the range; the headroom projection of Equation~\ref{eq:headroom} restores the recorded mass
without moving any cell outside it (Proposition~\ref{prop:headroom}). The projection has no
parameters.}
\label{fig:bounded}
\end{figure}

\clearpage

\section{Training, Initialization, and Measurement}
\label{sec:metrics}
\label{sec:setup}
Each training trajectory starts with a randomly drawn initial condition and a teacher
pre-seeding rollout of ten steps, except on Cahn--Hilliard where pre-seeding is disabled.
The teacher and emulator then evolve from the resulting common state over $K=12$ steps.
The optimizer receives the gradient of trajectory MSE through a \model{lax.scan} rollout,
with all epochs inside a single jitted scan and donated buffers. Evaluation uses a
distinct PRNG stream for held-out initial conditions and a four-times longer horizon.

\paragraph{One protocol, recorded per file.}
Every 2-D number in this paper
comes from one protocol: grid $48^2$, batch 32, 12 training steps, 48 evaluation steps,
1{,}200 epochs, He/Kaiming-normal initialization before ReLU, zero-initialized update
heads, 30 warmup epochs, and ten pre-seeding steps. The full matrix uses five seeds
$\{0,\dots,4\}$ and the three regime representatives use ten $\{0,\dots,9\}$. Each
results file records the JAX version, backend, device list, and the complete
configuration dataclass, so a table cannot be assembled from runs that differ in any of
these without the difference being visible.

Cahn--Hilliard's \model{preseed\_steps=0} is the one exception, and it is a documented
one: CH solutions coarsen rapidly from fresh initial conditions, so developed states
mismatch the evaluation distribution. Models within the CH comparison share the setting,
so the comparison remains internally controlled.

\begin{table}[t]
\centering
\caption{What the distillation target itself gets wrong, at the benchmark's own grid
($48^2$) and horizon. For the linear constant-coefficient members the periodic problem
is diagonal in Fourier space, so the continuum solution is available in closed form and
the error splits into the stencil's \emph{spatial} truncation and forward Euler's
\emph{temporal} error; for the rest the timestep is refined and the Cauchy differences
bound the teacher's own error, with the observed order of accuracy reported. Every
teacher converges. The last column is the reference solver's wall-clock per step
against the fastest emulator's.}
\label{tab:teacher}
\small
\resizebox{\linewidth}{!}{\begin{tabular}{llccccrr}
\toprule
Phenomenon & Estimate & Teacher $\relerr$ & Spatial & Temporal & Order & Solver (ms) & Emulator (ms) \\
\midrule
Heat             & closed form   & 0.00114 & 0.00181 & 0.00070 & --- & 1.64 & 3.86 \\
Wave             & closed form   & 0.00061 & 0.00048 & 0.00082 & --- & 3.79 & 4.23 \\
Advection--diff. & closed form   & 0.00472 & 0.00480 & 0.00103 & --- & 3.27 & 4.61 \\
Allen--Cahn      & $\Delta t$ refinement & 0.00017 & --- & --- & 0.97 & 3.63 & 6.44 \\
Gray--Scott      & $\Delta t$ refinement & 0.00497 & --- & --- & 1.03 & 1.70 & 1.26 \\
Shallow water    & $\Delta t$ refinement & $9.7\!\times\!10^{-7}$ & --- & --- & n/a$^{\dagger}$ & 4.92 & 4.72 \\
Cahn--Hilliard   & $\Delta t$ refinement & 0.0441 & --- & --- & 0.94 & 3.08 & 6.28 \\
FitzHugh--Nagumo & $\Delta t$ refinement & 0.00601 & --- & --- & 1.00 & 4.06 & 6.78 \\
Nagumo           & $\Delta t$ refinement & 0.00282 & --- & --- & 1.00 & 3.81 & 6.36 \\
Navier--Stokes   & $\Delta t$ refinement & 0.0251 & --- & --- & 1.54 & 1.15 & 1.28 \\
\bottomrule
\end{tabular}}
\\[2pt]
{\footnotesize $^{\dagger}$Shallow water's Cauchy differences reach machine precision,
so the order estimate is undefined; the teacher is fine, the estimator is not.}
\end{table}


\paragraph{Teacher timestep corrections.}
Two reference solvers required stability corrections that are recorded rather than
silently applied. Gray--Scott ships at $\Delta t=2.0$, above the explicit-diffusion bound
$\Delta x^2/(4D_u)=1.25$; experiments use $\Delta t=1.0$. Cahn--Hilliard is fourth order,
so its explicit bound is $\Delta t\le 2/(|\lambda_{\mathrm{lap}}|+\varepsilon^2
\lambda_{\mathrm{lap}}^2)=0.231$, and the original configuration used $\Delta t=0.5$. It
avoided visible blow-up only because the stepper clips the state each step; removing the
clip at that timestep reaches NaN, and timestep refinement gives an observed order of
accuracy of zero. Experiments therefore use $\Delta t=0.02$, where the scheme converges
at the expected first order (Table~\ref{tab:teacher}). The verbatim parameters are kept
in the registry so the migration gate still compares against the original. This is why
the teacher-convergence verdict in Table~\ref{tab:teacher} is reported per equation
rather than assumed.

\paragraph{Reported metric inventory.}
Relative $L^2$ error, MSE, RMSE, MAE, $L^\infty$, PSNR, an SSIM-style structural score,
high-frequency error fraction, relative error at $T/4,T/2,3T/4,T$, error-growth ratio,
per-channel mass drift and its drift curve, periodic-boundary residual, gradient energy,
parameter count, training wall time, inference time per step, and throughput. The primary
metric is per-initial-condition relative $L^2$, because pairing requires one value per
initial condition; the batch-reduced ratio of norms is also recorded and is a different
estimator of the same quantity, so the two do not print equal numbers.

Mass drift is $|\mass(\mathrm{prediction})-\mass(\mathrm{initial\ condition})|$, computed
\emph{per channel}: a lumped total would score a model that adds mass to one field and
removes it from another as perfectly conserving. For non-conservative targets a low drift
score indicates a wrong constraint rather than fidelity. Gradient energy is a diagnostic
proxy, not a proof of nonlinear stability, and a normalized high-frequency error fraction
need not decrease when total error decreases.

\begin{table}[t]
\centering
\caption{Diagnostic panel on heat (grid $48^2$, ten seeds). High-frequency error fraction
is the share of squared residual energy in the upper half of the spatial spectrum: it
describes \emph{where} the residual sits, not how large it is. Mass drift is the absolute
change in total heat at the end of the 48-step rollout. Error growth is the ratio of
relative error at the final step to that at one quarter of the horizon.}
\label{tab:heat20}
\small
\resizebox{\linewidth}{!}{\begin{tabular}{lcccccc}
\toprule
Model & $\relerr$ & PSNR (dB) & HF frac.\ & Mass drift & Err.\ growth & Train (s) \\
\midrule
FNO               & 0.001447 & 71.7 & 0.036 & 2.04 & 2.11 & 107.6 \\
MC-PI-NCA         & 0.003766 & 65.3 & 0.358 & $6.1\!\times\!10^{-5}$ & 5.47 & 18.2 \\
U-Net             & 0.004081 & 62.4 & 0.146 & 2.74 & 3.07 & 57.6 \\
ResNet-S          & 0.005559 & 60.1 & 0.063 & 7.83 & 2.74 & 20.7 \\
PI-NCA            & 0.006057 & 60.4 & 0.553 & $8.2\!\times\!10^{-5}$ & 5.52 & 12.9 \\
NCA               & 0.05614  & 46.1 & 0.331 & 65.3 & 22.87 & 15.3 \\
Identity (floor)  & 0.1889   & 29.3 & 0.000 & 0    & 3.70 & --- \\
\bottomrule
\end{tabular}}
\end{table}


\section{Additional Results and Resource Accounting}
\label{sec:additional-results}

\paragraph{Heat.}
Table~\ref{tab:heatms} gives the complete ten-seed ranking. Three observations that the
main text compresses. First, the standard NCA is an order of magnitude behind every other
learned model ($0.0561$ against $0.0074$ for the next worst) and its error-growth ratio
is $22.9\times$, against $5.5\times$ for PI-NCA and $2$--$3\times$ for the other models;
the deficit is rollout instability, not one-step accuracy. Second, three of the four
architectures between $0.0034$ and $0.0041$ are PI-NCA models (PI-NCA-Spec, PI-NCA-MS and
MC-PI-NCA), and the fourth is a U-Net with $265{,}104$ parameters, so within that accuracy
band parameter count is not predictive. Third, at $n=80$ every difference against the FNO
is Holm-significant.

\begin{table}[t]
\centering
\caption{Heat, ten seeds $\times$ 8 held-out initial conditions ($n=80$ paired
evaluations per model), grid $48^2$, 1{,}200 epochs. Intervals are 10{,}000-sample
percentile bootstraps on the per-initial-condition errors. Every model is compared with
the FNO by a Wilcoxon signed-rank test on the same initial conditions, Holm-corrected;
every difference is significant. The solver's own error on this task is $0.00114$, so the
FNO has nearly saturated the distillation target and the gap to the next rows should not
be over-interpreted.}
\label{tab:heatms}
\small
\begin{tabular}{lccrl}
\toprule
Model & $\relerr$ & 95\% CI & Params & vs.\ FNO \\
\midrule
FNO               & 0.001447 & [0.00137, 0.00154] & 592{,}897 & reference \\
MC-PI-NCA         & 0.003766 & [0.00310, 0.00453] & 9{,}936   & worse \\
U-Net             & 0.004081 & [0.00383, 0.00434] & 265{,}104 & worse \\
ResNet-S          & 0.005559 & [0.00520, 0.00594] & 5{,}364   & worse \\
FNO-S             & 0.005616 & [0.00531, 0.00592] & 8{,}433   & worse \\
PI-NCA (bounded)  & 0.005720 & [0.00505, 0.00644] & 4{,}576   & worse \\
PI-NCA            & 0.006057 & [0.00529, 0.00687] & 4{,}576   & worse \\
ResNet            & 0.006070 & [0.00542, 0.00677] & 74{,}336  & worse \\
U-Net-S           & 0.007381 & [0.00683, 0.00797] & 7{,}464   & worse \\
NCA               & 0.05614  & [0.0405, 0.0725]   & 6{,}784   & worse \\
Identity (floor)  & 0.1889   & [0.185, 0.193]     & 1         & worse \\
\bottomrule
\end{tabular}
\end{table}


\begin{table}[t]
\centering
\caption{Heat, grouped by parameter budget rather than ranked across budgets. Within a
class the comparison isolates the prior; across classes it reports what two orders of
magnitude more parameters actually bought. Inference is per rollout step at batch 8 on
the RTX 4090. Comparing the large spectral operator with a five-thousand-parameter
cellular automaton and calling the gap an architectural result would conflate capacity
with inductive bias, which is why the classes are separated.}
\label{tab:eff}
\small
\begin{tabular}{llcrr}
\toprule
Budget & Architecture & $\relerr$ & Params & Infer.\ (ms/step) \\
\midrule
\multirow{2}{*}{floor}
 & Identity                 & 0.1889   & 1       & 3.00 \\
\midrule
\multirow{8}{*}{small}
 & PI-NCA-MS    & 0.003699 & 5{,}520 & 1.79 \\
 & MC-PI-NCA            & 0.003766 & 9{,}936 & 1.32 \\
 & ResNet-S              & 0.005559 & 5{,}364 & 1.27 \\
 & FNO-S               & 0.005616 & 8{,}433 & 3.39 \\
 & PI-NCA-B       & 0.005720 & 4{,}576 & 1.25 \\
 & PI-NCA-MS-B & 0.005867 & 5{,}520 & 1.88 \\
 & PI-NCA                  & 0.006057 & 4{,}576 & 1.23 \\
 & U-Net-S                & 0.007381 & 7{,}464 & 1.58 \\
 & NCA               & 0.05614  & 6{,}784 & 1.48 \\
\midrule
medium & ResNet             & 0.006070 & 74{,}336 & 1.67 \\
\midrule
\multirow{3}{*}{large}
 & FNO                      & 0.001447 & 592{,}897 & 4.17 \\
 & PI-NCA-Spec      & 0.003405 & 134{,}225 & 2.77 \\
 & U-Net                     & 0.004081 & 265{,}104 & 1.80 \\
\bottomrule
\end{tabular}
\end{table}


\paragraph{Resource accounting.}
Table~\ref{tab:eff} groups by parameter budget. The within-class comparison is the one
that bears on inductive bias: at a matched small budget the conservative local models
lead the physics-free controls ($0.003699$ for PI-NCA-MS against $0.005559$ for
ResNet-S and $0.007381$ for U-Net-S).
Across classes, a $107\times$ increase in parameters buys a factor of $2.6$ in error.
Inference latency does not track parameter count: the $592{,}897$-parameter FNO costs
$4.17$\,ms per step against $1.23$\,ms for PI-NCA and $1.79$\,ms for PI-NCA-MS,
but the $265{,}104$-parameter U-Net costs $1.80$\,ms, because the FFT rather than the
weight count dominates. Accuracy, size, and latency must be read separately, and none of
these ratios establishes solver acceleration \citep{mcgreivy2024baselines}; the only
statement about that is Table~\ref{tab:matched}.

\paragraph{Ablation details.}
All four ablations run at the headline protocol with five seeds. A4 and A5 hold the
backbone fixed and change one component; A1 and A7 use an identical $5{,}520$-parameter
backbone and identical measured bounds, differing only in whether the output is clipped
and in how mass is restored afterwards. A4's parameter counts differ by $0.7\%$
($4{,}576$ against $4{,}544$), which is the cost of a two-channel flux head versus a
one-channel residual head and is far too small to explain a $4.6\times$ or $24\times$
effect.

A7 deserves one clarification. Its three variants are indistinguishable in accuracy
($0.3589$, $0.3614$, $0.3583$ on CH) and differ enormously in mass drift ($18.6$,
$1.07\times10^{-4}$, $5.40\times10^{-5}$). The correct reading is that the clip supplies
the stability and the projection supplies the conservation; they are separate mechanisms
that happen to be applied in sequence. The headroom projection of
Equation~\ref{eq:headroom} improves the drift by a further factor of two over the uniform
one and, unlike the uniform one, cannot move a clipped cell back outside the bound. We
present it as a guarantee obtained at no accuracy cost. The failure of the uniform
offset is provable (Appendix~\ref{sec:proofs}) and is exhibited by a constructed state in
the test suite; how often it occurs in trained rollouts at this scale is not measured.

\section{Evaluation Axes Not Measured at This Scale}
\label{sec:notmeasured}
Three axes present in the benchmark harness produced no results in this run, and we
report them as missing rather than filling them from earlier reduced-scale runs on
different hardware. Mixing backends within one study is precisely the error the controls
in this paper exist to prevent.

\begin{itemize}
\item \textbf{Out-of-distribution generalization.} The harness defines held-out axes
before training---amplitude, structure count, length scale, added broadband spectral
content, a random torus translation as a negative control, PDE coefficient, and
horizon---and evaluates a trained model on each. All three phenomena failed in this run
after 8--10 minutes each. The failure is scale-specific: the same command completes on a
smaller grid.
\item \textbf{Long-horizon stability with the guard disabled.} The accuracy tables clamp
rollouts to the teacher's measured physical range so that a diverging model cannot emit a
meaningless metric. That guard is correct for accuracy and wrong for stability, so the
stability study disables it and counts failures. All three phenomena failed in this run.
\item \textbf{Resolution transfer.} Skipped by the runner because a previous file was
present; the surviving file was produced on another backend and is therefore not used.
\end{itemize}

Because no stability numbers are available at this scale, the paper makes no long-horizon
robustness claim. The one stability-adjacent quantity that \emph{is} measured is the
error-growth ratio in Table~\ref{tab:heat20}, which is computed inside the guarded
accuracy protocol and is a weaker statement.

\section{Rank Stability Under Scale}
\label{sec:scaling-details}
\begin{table}[t]
\centering
\caption{Rank stability. Each axis is swept with everything else fixed at the main
operating point (grid $48^2$, rollout 12, 1{,}200 epochs, 48 evaluation steps, one seed),
for five models: NCA, PI-NCA, FNO, ResNet-S and the identity floor. Kendall's $\tau$
compares the ordering at the smallest and largest setting ($\tau=1$: unchanged). The
winner is reported at all three settings of each axis.}
\label{tab:scaling}
\small
\begin{tabularx}{\linewidth}{@{}llcX@{}}
\toprule
Phenomenon & Axis & $\tau$ & Winner at all three settings \\
\midrule
Heat & grid $24,32,48$           & 0.80 & FNO \\
     & rollout $4,8,12$          & 0.80 & FNO \\
     & epochs $300,600,1200$     & 0.80 & FNO \\
\midrule
Cahn--Hilliard & grid $24,32,48$       & 0.80 & PI-NCA \\
               & rollout $4,8,12$      & 0.60 & PI-NCA \\
               & epochs $300,600,1200$ & 0.60 & PI-NCA \\
\midrule
Navier--Stokes & grid $24,32,48$       & 1.00 & FNO \\
               & rollout $4,8,12$      & 0.00 & FNO \\
               & epochs $300,600,1200$ & 0.40 & FNO \\
\bottomrule
\end{tabularx}
\end{table}

Each sweep varies one axis with the others fixed at the headline operating point, using
six architectures spanning the unconstrained local, conservative local, global spectral,
and physics-free families, at one seed. Kendall's $\tau$ compares only the endpoints and
treats all rank inversions as equal, including inversions between models whose intervals
overlap; the per-setting errors in the released results distinguish the two cases.

On heat the grid sweep changes the winner: at $24^2$ PI-NCA-MS leads, and at $48^2$ the
FNO does. A local rule's reach per step is fixed in cells, so the number of steps it needs
to communicate across the domain grows with the grid, and its advantage on small domains
should not be extrapolated to large ones (Section~\ref{sec:scaling}). Cahn--Hilliard is
the least stable ordering, with all three axes changing it, consistent with its being the
stiffest system in the suite; at every one of its nine settings, however, the best of the
six swept architectures is PI-NCA or PI-NCA-MS.

\section{Matched Comparison Against a PINN and Against the Solver}
\label{sec:operators}
\begin{table}[t]
\centering
\caption{Matched comparison against a PINN, and against the solver. The task is fixed
rather than the model: given $K=32$ held-out initial conditions from the same
distribution, produce the solution of the heat equation at $T=48$ solver steps on a
$32^2$ grid, scored by the same metric against the same reference. The emulator trains
once and rolls out 32 times; the PINN trains once \emph{per initial condition}. The
paired test is on the same initial conditions ($p=4.7\times10^{-10}$, emulator better on
$100\%$ of them). Note the last row: at this grid the reference solver is both exact by
construction and cheaper per initial condition than either surrogate.}
\label{tab:matched}
\small
\begin{tabular}{lccrrrr}
\toprule
Method & $\relerr$ & 95\% CI & Params & Fixed (s) & Per IC (s) & Total (s) \\
\midrule
Emulator (\model{multiscale\_flux\_nca}) & 0.0153 & [0.0144, 0.0165] & 5{,}520 & 27.2 & 0.055 & 28.9 \\
PINN, one run per initial condition      & 0.1235 & [0.118, 0.129]   & 14{,}209 & 0 & 389.8 & 12{,}474.3 \\
Numerical solver (reference)             & 0 & --- & --- & 0 & 0.019 & 0.604 \\
\bottomrule
\end{tabular}
\end{table}

A PINN and an autoregressive emulator solve different problems, so the comparison fixes
the task rather than the model: produce the solution at $T$ for $K=32$ initial conditions
drawn from the same distribution, scored by the same metric against the same reference.
The emulator trains once and rolls out 32 times; the PINN trains once per initial
condition, at 6{,}000 iterations each. Costs are reported separately as a fixed component
and a per-instance component, so the amortization crossover is a derived quantity rather
than an assumption.

The emulator is better on every one of the 32 initial conditions while receiving
$431\times$ less total wall-clock, and the crossover is $K=1$. The solver row is included
because a surrogate must clear the solver, not only the other surrogate: at $32^2$ the
solver is roughly three times cheaper per initial condition than the emulator and exact
by construction. At the benchmark's own $48^2$ grid the gap narrows to near parity and
the emulator is faster on Gray--Scott and shallow water (Table~\ref{tab:teacher}). The
regime in which a learned surrogate is worthwhile therefore begins beyond the scales
measured here.

Standalone PINN, DeepONet, and steady Darcy baselines were also executed in this run
(284\,s, 3\,s, and 1{,}198\,s respectively) but write their results to the run log rather
than to a results file, so their numbers are not reported in this manuscript. The PINN
value that does appear, $0.1235$ in Table~\ref{tab:matched}, is from the matched
comparison and is the one that is well posed against the emulator track.

\section{Three-Dimensional Results}
\label{sec:three-details}
\begin{table}[t]
\centering
\caption{Three-dimensional results ($24^3$, 600 epochs, \textbf{single seed}, seven
architectures). The 3-D suite carries neither the physics-free CNN controls nor the
identity floor that the 2-D matrix uses, so these rows support a weaker statement than
Table~\ref{tab:regime2d}: they show that structural trends recur across dimensionality,
and nothing stronger. Two dimensionalities at one resolution and one seed are not
independent samples of the same population, so no ranking here transfers.}
\label{tab:regime3d}
\small
\begin{tabularx}{\linewidth}{@{}lllcX@{}}
\toprule
Phenomenon & Best in 3-D & Best in 2-D & 3-D $\relerr$ & Agreement \\
\midrule
Heat             & FNO       & FNO         & 0.00405 & same family \\
Advection--diff. & FNO       & FNO         & 0.00161 & same family \\
Allen--Cahn      & FNO       & PI-NCA-Spec & 0.00898 & both spectral \\
Nagumo           & NCA& U-Net        & 0.0273  & differs; both unconstrained \\
Gray--Scott      & NCA& U-Net        & 0.602   & differs; both unconstrained \\
FitzHugh--Nagumo & FNO       & NCA  & 0.152   & \textbf{differs} \\
\bottomrule
\end{tabularx}
\end{table}

The 3-D suite uses $24^3$ grids, 600 epochs, eight-step training and 24-step evaluation
horizons, a single seed, and seven architectures. It omits the physics-free ResNet and
U-Net controls and the identity floor, both of which the 2-D matrix uses, so it cannot
support the same kind of statement: with no floor, ``better than nothing'' is not
checkable, and with no physics-free control a conservative model leading a field of
conservative models would be uninformative.

Read with those caveats, the 2-D and 3-D winners agree in family on five of six
phenomena. The FNO leads heat, advection--diffusion, Allen--Cahn, and FitzHugh--Nagumo;
the unconstrained standard NCA leads Nagumo and Gray--Scott. FitzHugh--Nagumo is the one
disagreement, where the 2-D winner is the standard NCA and the 3-D winner is the FNO.
Conservative models remain clustered well behind on the reaction systems---on Nagumo all
five conservative variants sit at $0.197$ against the standard NCA's $0.0273$---which is the
same sign as the 2-D A4 ablation. Neither 3-D Navier--Stokes nor 3-D Cahn--Hilliard is
tested.

\section{Reproduction Record}
\label{sec:repro}
Every result in this paper was produced by one command on one machine. The run executed
52 stages over 22.2 hours on a single NVIDIA RTX 4090 with JAX 0.11.1, of which six
failed (Appendix~\ref{sec:notmeasured}) and one was skipped; the per-stage timings, exit
codes, and failures are recorded in the run manifest that accompanies the artifact.

The run begins with a correctness gate of 245 tests that must pass before any benchmark
executes. The gate asserts that the JAX ports equal verbatim PyTorch references to
tolerance, and it includes regression tests for two defects found during the migration:
a periodic Laplacian that differenced the wrong axes for the array layout the training
path actually uses, and the Cahn--Hilliard timestep instability described in
Appendix~\ref{sec:setup}. It also asserts properties the statistical claims rest on---that
the bootstrap interval covers at its nominal rate, that Holm correction is strictly more
conservative than uncorrected testing, that the closed-form solver references are correct,
and that the headroom projection is simultaneously mass-exact, bound-preserving, and
differentiable.

Two generated documents accompany the artifact and are not written by hand. A claims
audit derives the inventory of what was measured from the results files and evaluates
fifteen pre-registered claims against it, reporting each as supported, unsupported, or
not yet measured; for this run eleven are supported, one---rank stability---is
unsupported, and three are not yet measured. A bibliography check resolves every cited
arXiv identifier against the arXiv API. Both are regenerated at the end of every run, so
they describe the run that actually happened rather than an intended one.

\paragraph{A negative engineering result.} The CAX library for hardware-accelerated
cellular automata \citep{cax} was evaluated as a drop-in accelerator for the rollout and
measured to give no speedup on either CPU or GPU: both it and the \model{lax.scan} core
compile to the same XLA scan. It was never on the hot path and is not a dependency of the
released stack. We record this because a negative integration result is otherwise
invisible.

\paragraph{Remaining author checks.} The following should be completed before submission
and are not established by this manuscript: rerunning the three failed evaluation axes at
this scale; confirming the interpretation of mass drift on each reaction field;
supplying the standalone continuous-baseline numbers from the run log or rerunning them
with file output; and a personal review of the scientific interpretations and the
bibliography.
